% !TeX program = xelatex
\documentclass{UAmathtalk}
\renewcommand{\when}{12:00 p.m.}
\renewcommand{\where}{Hudson 0110}
\author{Jake Cordes}
\address{SUNY Albany}
\urladdr{https://www.albany.edu/math/people}
\title{Discussion on Persistent~Homology~Transforms}
\date{Friday, March 13, 2026}


\begin{document}

\maketitle

\begin{abstract}
\emph{Persistent Homology Transforms} (PHTs) are a generalization of applying persistent homology to level-set filtrations of height functions. 
Given a subset of Euclidean space $M$, the PHT on $M$ is a function from a sphere of appropriate dimension to the product of persistent diagrams, one for each dimension. I will be discussing basic definitions and results from the first paper on the subject \emph{"Persistent Homology Transform for Modeling Shapes and Surfaces"} by Katharine Turner, Sayan Mukherjee, and Doug M Boyer. The main focus of the talk will be the proof that the PHT is injective when restricted on shapes representable as finite simplicial complexes in 3-space. I will conclude with what work still needs to be done in the field and what I hope to accomplish in it.
\end{abstract}
\end{document}
